\section{Resource-Optimal First-Order Feature Geometry}
\label{sec:first-order-geometry}

Consider states
\[
z=(z_1,\ldots,z_m),
\qquad
z_i\overset{\mathrm{iid}}{\sim}\operatorname{Bernoulli}(p),
\]
with target distribution
\[
p_\star(x=i\mid z)
=
\frac{\exp(\alpha_i z_i)}
{\sum_{j=1}^m\exp(\alpha_j z_j)}.
\]
A valid target logit vector is therefore
\[
\ell_\star(z)=\alpha\odot z,
\]
up to addition of a constant shared across outputs.

The model assigns each complete state a free embedding
\[
E(z)\in\mathbb R^d
\]
and each output a shared decoder vector
\[
D(i)\in\mathbb R^d,
\]
with logits
\[
\ell_i(z)=E(z)^\top D(i).
\]
Because all states use the same \(d\)-dimensional decoder geometry, the full logit
matrix has rank at most \(d\).

We study the model's first-order dependence on the features. Define \(E_i\) as the
first-order coefficients of the least-squares fit
\[
E(z)\approx \bar E+\sum_i E_i z_i ,
\]
which under independent features equals the conditional-mean difference
\(\mathbb E[E(z)\mid z_i=1]-\mathbb E[E(z)\mid z_i=0]\) in population (the least-squares
form is what we compute: it removes the other features' chance sample correlations and
stays well defined under correlation). Collect the vectors \(E_i^\top\) as the rows of
\[
E^{(1)}\in\mathbb R^{m\times d}.
\]
Writing
\[
D=[D(1),\ldots,D(m)]\in\mathbb R^{d\times m},
\]
the first-order feature-to-logit map is
\begin{equation}
\boxed{
L^{(1)}=E^{(1)}D,
\qquad
\operatorname{rank}L^{(1)}\le d.
}
\label{eq:first-order-map}
\end{equation}

Our goal is to predict \(L^{(1)}\).

\subsection{Whitened first-order coordinates}

The factorization in Equation~\eqref{eq:first-order-map} is invariant under any
invertible change of hidden coordinates:
\[
E^{(1)}\mapsto E^{(1)}A,
\qquad
D\mapsto A^{-1}D.
\]
We use this freedom to whiten the first-order feature vectors:
\begin{equation}
\boxed{
E^{(1)\top}E^{(1)}=I_d.
}
\label{eq:whitened-gauge}
\end{equation}
This is a choice of coordinates, not an assumption about the model.

Define
\begin{equation}
\boxed{
d_i:=\|E_i\|^2.
}
\label{eq:dimension-allocation}
\end{equation}
Because the whitened vectors collectively satisfy
\[
\sum_iE_iE_i^\top=I_d,
\]
their squared norms obey
\begin{equation}
\boxed{
0\le d_i\le1,
\qquad
\sum_i d_i=d.
}
\label{eq:dimension-budget}
\end{equation}
Thus \(d_i\) is the effective amount of first-order dimensional capacity allocated
to feature \(i\).

\subsection{Convex allocation}

We first ask how the \(d\) available dimensions should be distributed among the
features, before choosing their detailed orientation.

The target logits are \(zW^\top\) with the logit map \(W=\operatorname{diag}(\alpha)\)
for independent features (Section~\ref{sec:modes} gives the general \(W\)). In the
feature-aligned fractional relaxation, retaining dimensional share \(d_i\) scales
feature \(i\)'s contribution by \(d_i\):
\[
q_{\mathbf d}(\cdot\mid z)
=
\operatorname{softmax}\!\big((z\odot\mathbf d)\,W^\top\big),
\qquad\text{i.e. } q_{\mathbf d}(i\mid z)\propto\exp(\alpha_i d_i z_i)\ \text{when } W=\operatorname{diag}(\alpha).
\]
We predict the allocation by
\begin{equation}
\boxed{
\mathbf d^\star
=
\arg\min_{\substack{0\le d_i\le1\\ \sum_i d_i=d}}
\mathbb E_z\,
\operatorname{CE}
\left[
p_\star(\cdot\mid z),
q_{\mathbf d}(\cdot\mid z)
\right].
}
\label{eq:allocation-program}
\end{equation}
This is a convex optimization because the logits are linear in \(\mathbf d\) and
cross-entropy is convex in the logits.

Define the marginal predictive utility of allocating more dimension to feature
\(i\) by
\[
u_i(\mathbf d)
=
-
\frac{\partial\operatorname{CE}(\mathbf d)}
{\partial d_i}.
\]
Differentiating gives
\begin{equation}
\boxed{
u_i(\mathbf d)
=
\alpha_i p
\left[
\mathbb E[p_{\star i}\mid z_i=1]
-
\mathbb E[q_{\mathbf d,i}\mid z_i=1]
\right].
}
\label{eq:marginal-utility}
\end{equation}

Let \(\delta\) be the shadow price of one unit of dimensional capacity. The KKT
conditions for Equation~\eqref{eq:allocation-program} are
\begin{equation}
\boxed{
\begin{cases}
u_i(\mathbf d^\star)\le \delta,
&d_i^\star=0,
\\[1mm]
u_i(\mathbf d^\star)=\delta,
&0<d_i^\star<1,
\\[1mm]
u_i(\mathbf d^\star)\ge \delta,
&d_i^\star=1.
\end{cases}
}
\label{eq:water-filling}
\end{equation}
Thus features of different importance receive different amounts of capacity until
the next infinitesimal unit of dimension has the same predictive value wherever it
is used.

\subsection{Euclidean realization and superposition}

After solving Equation~\eqref{eq:allocation-program}, choose first-order feature
vectors satisfying
\begin{equation}
\boxed{
\|E_i\|^2=d_i^\star,
\qquad
\sum_iE_iE_i^\top=I_d.
}
\label{eq:tight-frame-realization}
\end{equation}
Such vectors form a Parseval tight frame. A realization exists whenever
Equation~\eqref{eq:dimension-budget} holds, although it is generally not unique.

For every such realization,
\begin{equation}
\boxed{
\sum_{j\neq i}
\langle E_i,E_j\rangle^2
=
d_i^\star(1-d_i^\star).
}
\label{eq:forced-overlap}
\end{equation}
Hence fractional dimensional allocation necessarily implies superposition:
\[
0<d_i^\star<1
\quad\Longrightarrow\quad
\langle E_i,E_j\rangle\neq0
\quad\text{for at least one }j\neq i.
\]
The allocation predicts the total squared overlap of each feature, but not its
individual overlap partners or signs.

\subsection{Predicting the first-order logit map}

For a chosen compatible feature geometry \(E^{(1)}\), we solve directly for the
shared decoder together with a mean-embedding intercept
\(e_0\in\mathbb R^d\):
\begin{equation}
\boxed{
\left(D^\star,e_0^\star\right)(E^{(1)})
=
\arg\min_{D,\,e_0}
\mathbb E_z\,
\operatorname{CE}
\left[
p_\star(\cdot\mid z),
\operatorname{softmax}
\left(
\left(z^\top E^{(1)}+e_0^\top\right)D
\right)
\right].
}
\label{eq:decoder-program}
\end{equation}
The intercept is the \(\bar E\) coordinate of the state decomposition: it
contributes the state-independent logit row \(e_0^\top D\), carried through the
decoder — the model class has no per-output bias. We fix the softmax gauge by
requiring
\[
D\mathbf 1=0
\]
in the analysis; the implementation fixes the same invariance by row-centring every
model's logit table over states before factorising it, which is the same gauge applied
per state.
For fixed \(E^{(1)}\), Equation~\eqref{eq:decoder-program} is convex in \(D\) at
fixed \(e_0\) and in \(e_0\) at fixed \(D\); the pair is solved jointly.

The predicted first-order map is then
\begin{equation}
\boxed{
L^{(1)}
=
E^{(1)}D^\star(E^{(1)}).
}
\label{eq:predicted-first-order-map}
\end{equation}

Because Equation~\eqref{eq:tight-frame-realization} does not determine a unique
orientation, the theory predicts a family of first-order maps, one for each
compatible tight-frame realization. In the iid setting, the stable predictions are
the allocation \(d_i^\star\), the unavoidable overlap in
Equation~\eqref{eq:forced-overlap}, and statistics of the resulting maps; individual
pairwise overlaps may remain underdetermined.

Finally, write
\[
L^{(1)}=U\Sigma V^\top.
\]
The matrices
\[
UU^\top,\qquad VV^\top,\qquad \Sigma
\]
give the invariant feature-side geometry, decoder-side geometry, and predictive-mode
strengths.

The map alone does not fix raw feature norms: for any invertible \(A\),
\((EA)(A^{-1}D)=ED\), and the loss is exactly flat along that direction. Plain
gradient descent from small initialization pins it, since \(E^\top E-D^\top D\) is
conserved, so the solution is balanced; the trained models here use Adam, whose
per-coordinate normalization does not conserve it (the state table \(E\) ends several
times larger than \(D\); retrained with SGD the two agree to \(0.4\%\)), and the
predicted model is stored as the completion returns it (\(E^\star\) a Parseval frame,
the scale in \(D\)). Models are therefore compared through quantities that do not
depend on \(A\), and every model's first-order vectors are read by one estimator, the
least-squares fit \(h(z)\approx\bar h+\sum_iE_iz_i\) of its saved representation.

\emph{Allocation and Grams} are read in the whitened coordinates of
Equation~\eqref{eq:whitened-gauge}, \(\tilde E=E^{(1)}(E^{(1)\top}E^{(1)})^{-1/2}\).
The leverage \(d_i=\|\tilde E_i\|^2\) is the diagonal of the projector
\begin{equation}
\boxed{
\tilde E\tilde E^\top
=
E^{(1)}\big(E^{(1)\top}E^{(1)}\big)^{+}E^{(1)\top},
}
\label{eq:whitened-grams}
\end{equation}
and that projector is the Gram we display: the allocation on its diagonal, the signed
overlaps of the frame off it, no scale left to choose. The readout is read the same way,
\(D^\top(DD^\top)^{+}D\). Across seeds we compare the shared frame rather than the average
Gram: each seed's whitened frame is carried onto the first predicted frame by orthogonal
Procrustes over all rows in \(\mathbb R^d\), the frames are averaged, and the Gram of the
mean frame is shown. The mean of per-seed Grams keeps every seed's private wobble as
haze; the Gram of the mean frame keeps what the seeds agree on. Placement is the
correlation of two such Grams' off-diagonal entries.

\emph{Sizes} need a scale, and the loss offers none, so one gauge is imposed on every
model's row-centred logit table over states (for the completed model
\(\Lambda=[Z,\ \phi(Z),\ \mathbf 1]\,[E^\star;E_2;e_0^\top]\,D\)), written
\(\Lambda=U\Sigma V^\top\), truncated to rank \(d\) and split evenly:
\begin{equation}
\boxed{
E_{\mathrm{bal}}
=
N^{1/4}\,U\Sigma^{1/2},
\qquad
D_{\mathrm{bal}}
=
N^{-1/4}\,\Sigma^{1/2}V^\top,
}
\label{eq:balanced-factorization}
\end{equation}
the factor \(N^{1/4}\) putting the state table on a per-state scale,
\(E_{\mathrm{bal}}^\top E_{\mathrm{bal}}/N=\operatorname{diag}(\sigma_k/\sqrt N)\).
Sizes are the \(\|E_i\|^2\) of the first-order vectors of \(E_{\mathrm{bal}}\); the Bayes
table itself reads about \(\alpha_i/\sqrt{p(1-p)}\) on this scale, so amplification is a
ratio to that reference rather than to \(\alpha_i\).

\subsection{Correlated features: the predictive modes}
\label{sec:modes}

Every stage above evaluates its objective on the data, so nothing in it assumes
independent features; what does depend on the coordinates is the relaxation itself.
The fractional predictor scales each coordinate's logit by its own leverage, and the
Fantope constrains an \emph{orthogonal} projector, so ``fraction represented'' is
measured in whichever basis the inputs are written. With correlated features that
basis matters: a shared component of two features is one direction, not two, and the
raw coordinates charge for it twice (a duplicated feature is the extreme case).

We therefore write the inputs in their predictive modes,
\begin{equation}
u = z\,C^{-1/2},\qquad C=\operatorname{Cov}(z),
\label{eq:modes}
\end{equation}
the symmetric root so that mode \(k\) stays closest to feature \(k\), and \(u\) has
identity covariance. The Bayes logits \(\alpha_i z_i\) become \(u\,W^\top\) with the
logit map
\begin{equation}
W=\operatorname{diag}(\alpha)\,C^{1/2}
\label{eq:mode-map}
\end{equation}
(outcomes by modes): outcome \(i\) reads a mixture of modes, which is the only
structural change. Water filling, pressure, Fantope program, frame, and completion
run unchanged on \((W, u)\); for a diagonal map they are the formulas of the previous
sections, and for independent features \(C\propto I\) so the modes are the features.
The trained networks are measured on the same coordinates (least squares of the
representation on \([1, u]\)); feature-space vectors follow as \(E = C^{-1/2}\tilde E\).

The choice is a modelling decision about the relaxation, not a theorem, and it is
what makes correlated features work: across the sweep (\(\rho\in\{0.3,0.6\}\), every
\(p\), \(d\)), for ladder slopes \(\beta\le 0.5\) the theory on the modes reaches the additive
class's own optimum to within 0--20\% in CE with leverage correlations 0.92--0.98, where
the raw coordinates charge twice for a shared component and fall short. The one regime it
misses is the steep ladder \(\beta=1\) at \(d\ge 6\), by a factor growing with \(\rho\) and \(d\)
(1.5 to 3.9); there the class optimum is hard and the refund assumption below fails
(\texttt{notes/superposition\_results.md}). On a mutually exclusive pair inside independent
features both coordinates predict the antipodal arrangement, the raw ones with the lower CE.

\subsection{Orientation selection: which overlaps, with which signs}
\label{sec:orientation}

The allocation fixes each feature's total overlap; the tight-frame manifold leaves
its placement and signs free. Trained networks are not uniform on that manifold:
the frame their seeds share keeps negative overlaps toward the head. The selection is
the odd term of the expansion whose even term prices overlap magnitudes. At the
fractional predictor \(q_{\mathbf d^\star}\), the first-order CE pressure on a signed
overlap \(\langle E_i,E_j\rangle\) is
\begin{equation}
\boxed{
a_{ij}
=
\alpha_j\,\mathbb E\!\left[(z_i-p)\,(q_j-p_{\star j})\right]
+
\alpha_i\,\mathbb E\!\left[(z_j-p)\,(q_i-p_{\star i})\right],
}
\label{eq:signed-pressure}
\end{equation}
so that to first order \(\Delta\mathrm{CE}\approx\sum_{i<j}a_{ij}\langle E_i,E_j\rangle\)
and \(a_{ij}>0\) prefers \(\langle E_i,E_j\rangle<0\). The mechanism is the flatness
of the compressed predictor: when \(z_i=1\) the competitors \(j\) are over-predicted,
\(q_j>p_{\star j}\), and a negative overlap lets \(i\) suppress \(j\). This term is
invisible to the convex program (every term there is even in the off-diagonals) and
to the uniform frame ensemble (whose mean cancels it), which is why neither predicts
signs.

\paragraph{Prediction rule.}
The predicted orientation follows the pressure as far as the Fantope allows, with the
diagonal free:
\begin{equation}
\boxed{
P^\star
=
\arg\min_{P\in\mathcal F_d}\;
C(\operatorname{diag}P)+\sum_{i<j}a_{ij}\,P_{ij},
\qquad
\mathcal F_d=\{0\preceq P\preceq I,\ \operatorname{tr}P=d\},
}
\label{eq:pressure-frame}
\end{equation}
the exact fractional-predictor CE of the diagonal plus the linearised interference of
the off-diagonals: one convex program, no decoder solve. Its optimum is not a
projector; \(\eta=\operatorname{tr}(P^\star-P^{\star2})\) measures the gap, and the
frame \(E^\star\) is the rank-\(d\) projector with the same diagonal as \(P^\star\)
nearest to it (alternating projections from \(P^\star\)), so the fractional leverages
and the soft tail of the next paragraph survive realization while \(\eta\) is paid
there. (The same program with the exact tied-map CE in place of the linearisation is
also convex, but its optimum leaves the projectors, \(\eta\approx1\), and loses the
orientation once realized; the linear off-diagonal term is what keeps the solution
near the projectors.) The rule is a statement about the optimum of the full class, not only
about its gradient: with the frame held fixed the class's loss is nearly flat across
placements (decoder and pairs compensate), but when frame, decoder and pairs are
optimized together the pressure-following placement wins clearly, and gradient
descent started from the opposite sign pattern abandons it and ends at the pressure
pattern at the same loss as every other start. Selecting frames by the
additive class's cross-entropy fails because that class cannot repair the
interference the pressure-following frames create. Evidence: \texttt{notes/superposition\_results.md}.

\paragraph{The soft tail.}
The same pressure gives every feature that water filling drops a second lever. A
feature with \(d_j=0\) cannot raise its own logit, but a small vector anti-aligned
with the represented features lowers their logits when \(z_j=1\), which raises
\(q_j\) relatively. The benefit is linear in the overlaps and, by Cauchy--Schwarz,
at most \(\|a_{\cdot j}\|\sqrt{d_j}\) for leverage \(d_j\), while the dimension costs
\(\delta d_j\); the optimum is small but never zero,
\begin{equation}
\boxed{
d_j=\frac{\|a_{\cdot j}\|^2}{4\delta^2}
\qquad\text{for features below the water-filling threshold,}
}
\label{eq:soft-tail}
\end{equation}
with \(\|a_{\cdot j}\|\) the pressure norm onto the represented set. Trained tail
leverage follows this law with no fitted constant where features are sparse and the
rank low (\(p\le0.15\), \(d\le6\): log--log slope \(\approx2\)); elsewhere the slope is
between 0.4 and 1.8, so the law is a sparse, low-rank statement. The soft tail is where
the trained interference beyond the predicted support lives. At \(d=12\) the dropped
features' leverage is flat and uncorrelated with the pressure — a converged property
that the water level, twenty times smaller there, no longer sets
(\texttt{notes/superposition\_results.md}).

\paragraph{The hierarchy, and what is specific to softmax.}
Three layers, each tested separately:
\begin{itemize}
\item \textbf{Magnitude}: water filling plus tight-frame geometry
(Equation~\eqref{eq:forced-overlap}).
\item \textbf{Coarse sign}: generic to compression. The quadratic-loss analogue of
the pressure makes the same sign calls: any loss facing a too-flat compressed
predictor pushes overlaps negative.
\item \textbf{Graded placement}: specific to the softmax (Fisher) weighting of those
compression errors. Only the softmax \(a_{ij}\) correlates with the trained Gram's
values; trained frames sit at the low-\(A\) tail of the compatible ensemble,
\(A=\sum_{i<j}a_{ij}\langle E_i,E_j\rangle\); and under a global logit scale
\(\alpha\mapsto s\alpha\) the softmax and quadratic pressures converge as \(s\to0\)
and separate as saturation grows.
\end{itemize}

\paragraph{Limit.}
At \(d=12\) the trained placement is a structured, mutually negative arrangement among
the fractional tail (a packing of near-equal features, \(\alpha\approx2\)--\(3.6\)). The
pressure frame reproduces about half of it on the shared frame (placement \(0.46\),
against \(0.67/0.73\) at \(d=3/6\)), and so do the pairs model and, at this \(d\) only,
the additive model. The remainder is not carried by pairs of non-represented features,
is not a head-versus-tail question, and is not resolved by the second-order term; which
tail partners pair off is decided by something the pressure order does not contain
(\texttt{notes/superposition\_results.md}).

\subsection{Stage two: repair at the predicted frame}
\label{sec:repair}

The first-order theory stops at \(L^{(1)}=E^{\star}D^\star\). The network's remaining
context dependence \(E_{\ge2}(z)\) is dominated by feature pairs, so we predict it by
the additive-plus-pairs class at the fixed frame \(E^\star\): code
\[
E(z)=z^\top E^\star+\phi(z)^\top E_2+e_0^\top,
\qquad
\phi_{ij}(z)=\text{the product }(z_i-p_i)(z_j-p_j)\text{ residualised on }[1,z]\text{ and standardised},
\]
for \(i<j\) among the top 18 features (153 pairs), read through one decoder \(D\). The
residualisation makes the pair block exactly orthogonal to the first-order dictionary in
the sample, so \(E^{(1)}\) is untouched by construction and the split remains well
defined when the features are correlated. The problem is convex
in \(E_2\) at fixed \((D,e_0)\) and in \((D,e_0)\) at fixed code; we alternate the
two solves from the additive decoder (block-coordinate descent to a stationary point,
not a certified optimum — the one non-certified step in the pipeline).

\paragraph{What it predicts.}
Two things the first-order theory could not: most of the cross-entropy gap to the
trained network, and the head amplification. The mechanism is the one anticipated in
the appendix: once the pair vectors absorb the interference the frame's overlaps
create, the decoder can read the head harder without paying for it on co-active
states. The completion has to be solved: a one-step quadratic in decoder resizing and
pair vectors at the tied baseline is outside its radius, and the pair sector is
conditional on the decoder as much as on the frame (numbers in
\texttt{notes/superposition\_results.md}).

\subsection{Status: validated, failed, and open}

All comparisons use five runs per budget \(d\in\{3,6,12\}\), with task and training
seeds coupled so that averages span task draws as well as initializations. The
experiments and numbers are recorded in \texttt{notes/superposition\_results.md}.

\paragraph{Validated.}
\begin{itemize}
\item \textbf{Allocation.} Water filling matches trained leverage at correlation
\(0.83\)--\(1.00\) over the whole sweep (\(p\), \(\beta\), \(\rho\), \(d\)), including
the fractional head where the regime is soft and the position of the support edge
(Figure~1 of the analysis notebook shows the full profiles, no threshold applied).
\item \textbf{Overlap magnitudes.} Per-row off-diagonal energy is fixed by
Equation~\eqref{eq:forced-overlap}; realized frames satisfy it exactly.
\item \textbf{Repair.} At \(E^\star\), the pair block closes most of the CE gap to
trained and reproduces the head amplification (Section~\ref{sec:repair}).
\item \textbf{Orientation.} The frame of Equation~\eqref{eq:pressure-frame} reproduces
the trained placement (off-diagonal correlation of the whitened Grams of the shared
frame, top 16 features) at \(0.67/0.73/0.46\) for \(d=3/6/12\) at the reference cell,
above the agreement of one trained seed with the rest (\(0.54/0.44/0.26\)) and level
with the pairs model (\(0.65/0.69/0.50\)); orientation is a joint-optimum property
(gradient descent from the anti-pressure frame returns to it) and is optimizer
independent (Adam and SGD agree within seed noise).
\end{itemize}

\paragraph{Known failures.}
\begin{itemize}
\item \textbf{Head amplification beyond pairs.} Stage two recovers the head at
\(d=3,6\) (Section~\ref{sec:repair}) but leaves the third feature at \(d=3\) short and,
at \(d=12\), the whole represented block (head \(16.5\) against \(20.1\) on the common
scale); the first-order theory alone under-predicts the whole block at every \(d\).
\item \textbf{Placement at \(d=12\).} The pressure frame reproduces the placement at
\(d=3,6\) and about half of it at \(d=12\), where the reference models do no better
(Section~\ref{sec:orientation}).
\item \textbf{The refund assumption.} Stage one minimises the tied predictor's CE
\(C(d)\), which counts the benefit of leverage and none of its cost: realising a fractional
diagonal forces the off-diagonal energy of Equation~\eqref{eq:app-forced-overlap}, whose
interference cost is the \(B(h)=\sum_i h_i(1-h_i)\) term of the two-term theory. Dropping
the projector constraint deletes \(B\), and that is correct exactly when the completion
refunds it. It does at the reference cell (free pairs return 78--84\% of the \(B\)-slope,
so a priced allocation there is worse, not better), and it fails where the ladder is steep
or the features correlated: there the additive class's own optimum is hard (about \(d\)
features at full leverage), the theory misses it by a factor of four in CE, and the
fractional allocation is decided by ties among near-equal tail importances. The
refund fails because the Fisher weight \(q_i(1-q_i)\) of the interference cost varies with
the state once the softmax is peaked, pushing the cost into Walsh levels the pair block
cannot reach. The candidate repair is a \emph{residual} price, the \(B\)-cost minus its
projection onto \(\operatorname{span}\{1,z,\phi(z)\}\), which is computable at the
fractional predictor and reduces to plain water filling where pairs refund
(\texttt{notes/superposition\_results.md}).
\item \textbf{Tail edges.} Trained allocations hold small leverage past the
water-filling cutoff. The free diagonal of Equation~\eqref{eq:pressure-frame} produces
such a tail, but a smaller one than trained at \(d=3,6\); the closed form of
Equation~\eqref{eq:soft-tail} matches the trained tail there and is not part of the
predicted allocation.
\end{itemize}

\paragraph{Open.}
The \(d=12\) placement mechanism; the residual of stage two (wider pair sets, triples);
the extended dictionary Gram (features
and pairs) per realization; the distortion--rate frontier
\(\Delta\) versus \(d\) with slope \(\delta\), connecting the allocation to
rate--distortion and information-bottleneck accounts in which rate is dimension.

\appendix

\section{Derivation of the First-Order Geometry}
\label{app:first-order-geometry}

\subsection{The rank-limited predictive model}

Let
\[
z=(z_1,\ldots,z_m),
\qquad
z_i\overset{\mathrm{iid}}{\sim}\operatorname{Bernoulli}(p),
\]
and
\[
p_\star(x=i\mid z)
=
\frac{\exp(\alpha_i z_i)}
{\sum_j\exp(\alpha_jz_j)}.
\]
A valid target logit vector is
\[
\ell_\star(z)=\alpha\odot z.
\]
Logits are defined only up to addition of a scalar shared across outputs.

The model assigns each state \(z\) an embedding
\[
E(z)\in\mathbb R^d
\]
and each output \(i\) a shared decoder vector
\[
D(i)\in\mathbb R^d.
\]
There is no decoder bias. The logits are
\[
\ell_i(z)=E(z)^\top D(i).
\]

Stacking all state embeddings and decoder vectors gives
\[
\widehat L=\mathsf E D,
\]
and therefore
\[
\operatorname{rank}\widehat L\le d.
\]
The state embeddings are free, but every state must be decoded through the same
\(d\)-dimensional output geometry.

\subsection{The first-order feature vectors}

In general, and in the implementation, \(E_i\) are the first-order coefficients of the
least-squares fit of \(E(z)\) on \([1,z]\). For independent features these coincide in
population with
\begin{equation}
E_i
=
\mathbb E[E(z)\mid z_i=1]
-
\mathbb E[E(z)\mid z_i=0],
\label{eq:app-feature-vector}
\end{equation}
which is the form the derivation below uses.
Let
\[
\bar E=\mathbb E[E(z)]
\]
and define the residual
\begin{equation}
E_{\ge2}(z)
=
E(z)
-
\bar E
-
\sum_i(z_i-p)E_i.
\label{eq:app-residual}
\end{equation}

For each \(i\),
\[
\mathbb E[(z_i-p)E(z)]
=
p(1-p)
\left(
\mathbb E[E(z)\mid z_i=1]
-
\mathbb E[E(z)\mid z_i=0]
\right)
=
p(1-p)E_i.
\]
Independence also gives
\[
\mathbb E[(z_i-p)(z_j-p)]=0
\qquad
i\neq j,
\]
and
\[
\mathbb E[(z_i-p)^2]=p(1-p).
\]
Therefore
\begin{equation}
\mathbb E[(z_i-p)E_{\ge2}(z)]=0
\qquad
\forall i.
\label{eq:app-residual-orthogonal}
\end{equation}
Under correlated features this step fails, which is why the implementation defines the
split by least squares on \([1,z]\) and residualises the pair block on the same
dictionary: orthogonality then holds in the sample by construction rather than by
independence.

Thus
\[
E(z)
=
\bar E
+
\sum_i(z_i-p)E_i
+
E_{\ge2}(z)
\]
is the decomposition of the state embedding into its first-order feature dependence
and the remaining context dependence.

Collect the \(E_i^\top\) as rows of
\[
E^{(1)}
=
\begin{bmatrix}
E_1^\top\\
\vdots\\
E_m^\top
\end{bmatrix}
\in\mathbb R^{m\times d},
\]
and write
\[
D=[D(1),\ldots,D(m)]
\in\mathbb R^{d\times m}.
\]
The first-order feature-to-logit map is
\begin{equation}
\boxed{
L^{(1)}=E^{(1)}D.
}
\label{eq:app-first-order-map}
\end{equation}
Consequently,
\begin{equation}
\boxed{
\operatorname{rank}L^{(1)}\le d.
}
\label{eq:app-first-order-rank}
\end{equation}

\subsection{Whitening as a hidden-coordinate choice}

The factorization of \(L^{(1)}\) is not unique. For every invertible
\(A\in\mathbb R^{d\times d}\),
\[
E^{(1)}D
=
(E^{(1)}A)(A^{-1}D).
\]
The same transformation may be applied to every state embedding, so it changes only
the hidden coordinates and not the model's logits.

Assume that \(E^{(1)}\) has rank \(d\). Starting from an arbitrary factorization,
apply
\begin{equation}
E^{(1)}
\mapsto
E^{(1)}
\left(E^{(1)\top}E^{(1)}\right)^{-1/2},
\label{eq:app-whiten-E}
\end{equation}
and
\begin{equation}
D
\mapsto
\left(E^{(1)\top}E^{(1)}\right)^{1/2}D.
\label{eq:app-whiten-D}
\end{equation}
Their product is unchanged, while the transformed feature matrix satisfies
\begin{equation}
\boxed{
E^{(1)\top}E^{(1)}=I_d.
}
\label{eq:app-whitened-gauge}
\end{equation}

We use the same symbols \(E^{(1)}\) and \(D\) after this coordinate change.

Equation~\eqref{eq:app-whitened-gauge} is equivalently
\begin{equation}
\boxed{
\sum_iE_iE_i^\top=I_d.
}
\label{eq:app-tight-frame}
\end{equation}
Thus the whitened first-order vectors form a Parseval tight frame.

Define
\begin{equation}
\boxed{
d_i=\|E_i\|^2.
}
\label{eq:app-di}
\end{equation}
Summing gives
\[
\sum_i d_i
=
\sum_i\operatorname{tr}(E_iE_i^\top)
=
\operatorname{tr}
\left(
\sum_iE_iE_i^\top
\right)
=
d.
\]
Moreover, the matrix
\[
E^{(1)}E^{(1)\top}
\]
is an orthogonal projector, so its diagonal entries satisfy
\[
0\le d_i\le1.
\]
Therefore
\begin{equation}
\boxed{
0\le d_i\le1,
\qquad
\sum_i d_i=d.
}
\label{eq:app-dimensional-budget}
\end{equation}

If the first-order component has rank \(r<d\), the same construction applies on its
support and the total budget is \(r\).

\subsection{The fractional allocation problem}

The hard geometry requires choosing all vectors \(E_i\) subject to
Equation~\eqref{eq:app-tight-frame}. We first solve the simpler question of how their
total squared norm should be divided among features.

The target logits are \(zW^\top\), \(W=\operatorname{diag}(\alpha)\) for independent
features and \(W=\operatorname{diag}(\alpha)C^{1/2}\) on the predictive modes. In the
feature-aligned fractional relaxation, retaining dimensional share \(d_i\) scales
coordinate \(i\)'s contribution by \(d_i\):
\begin{equation}
q_{\mathbf d}(\cdot\mid z)
=
\operatorname{softmax}\!\big((z\odot\mathbf d)W^\top\big),
\qquad
q_{\mathbf d}(i\mid z)
=
\frac{\exp(\alpha_i d_i z_i)}
{\sum_j\exp(\alpha_jd_jz_j)}
\ \text{ for } W=\operatorname{diag}(\alpha).
\label{eq:app-fractional-predictor}
\end{equation}

Define
\begin{equation}
\mathcal C(\mathbf d)
=
\mathbb E_z
\left[
-\sum_i
p_\star(i\mid z)
\log q_{\mathbf d}(i\mid z)
\right].
\label{eq:app-allocation-loss}
\end{equation}
The allocation is
\begin{equation}
\boxed{
\mathbf d^\star
=
\arg\min_{\substack{0\le d_i\le1\\ \sum_i d_i=d}}
\mathcal C(\mathbf d).
}
\label{eq:app-allocation-program}
\end{equation}

The logits in Equation~\eqref{eq:app-fractional-predictor} are linear in
\(\mathbf d\), cross-entropy is convex in the logits, and the constraints are convex.
Hence Equation~\eqref{eq:app-allocation-program} is a convex optimization problem.

This feature-aligned fractional program is the central allocation approximation. It
predicts how much first-order dimensional capacity each feature receives before the
detailed tight-frame orientation is chosen.

\subsection{Marginal predictive utility}

For one state \(z\), softmax cross-entropy satisfies
\[
\frac{\partial\operatorname{CE}_z}{\partial\ell_i}
=
q_{\mathbf d,i}(z)-p_{\star i}(z).
\]
The relaxed logit obeys
\[
\frac{\partial\ell_i(z)}{\partial d_i}
=
\alpha_i z_i.
\]
The chain rule therefore gives
\[
\frac{\partial\mathcal C}{\partial d_i}
=
\alpha_i
\mathbb E_z
\left[
z_i
\left(
q_{\mathbf d,i}-p_{\star i}
\right)
\right].
\]

Since \(z_i\in\{0,1\}\) and \(\Pr(z_i=1)=p\),
\[
\mathbb E[z_if(z)]
=
p\,\mathbb E[f(z)\mid z_i=1].
\]
Thus
\begin{equation}
\frac{\partial\mathcal C}{\partial d_i}
=
\alpha_i p
\left[
\mathbb E[q_{\mathbf d,i}\mid z_i=1]
-
\mathbb E[p_{\star i}\mid z_i=1]
\right].
\label{eq:app-allocation-gradient}
\end{equation}

Define the marginal predictive utility
\begin{equation}
\boxed{
u_i(\mathbf d)
=
-
\frac{\partial\mathcal C}{\partial d_i}.
}
\label{eq:app-utility-definition}
\end{equation}
Then
\begin{equation}
\boxed{
u_i(\mathbf d)
=
\alpha_i p
\left[
\mathbb E[p_{\star i}\mid z_i=1]
-
\mathbb E[q_{\mathbf d,i}\mid z_i=1]
\right].
}
\label{eq:app-utility}
\end{equation}

This is the expected reduction in predictive loss obtained by allocating an
additional infinitesimal amount of dimension to feature \(i\).

\subsection{Water filling}

Introduce a multiplier \(\delta\) for
\[
\sum_i d_i=d,
\]
and nonnegative multipliers \(\lambda_i,\mu_i\) for the lower and upper bounds. The
Lagrangian is
\[
\mathscr L
=
\mathcal C(\mathbf d)
+
\delta\left(\sum_i d_i-d\right)
-
\sum_i\lambda_i d_i
+
\sum_i\mu_i(d_i-1).
\]

Stationarity gives
\[
\frac{\partial\mathcal C}{\partial d_i}
+\delta-\lambda_i+\mu_i=0.
\]
Using
\[
u_i=-\frac{\partial\mathcal C}{\partial d_i},
\]
together with complementary slackness, gives
\begin{equation}
\boxed{
\begin{cases}
u_i(\mathbf d^\star)\le \delta,
&d_i^\star=0,
\\[1mm]
u_i(\mathbf d^\star)=\delta,
&0<d_i^\star<1,
\\[1mm]
u_i(\mathbf d^\star)\ge \delta,
&d_i^\star=1.
\end{cases}
}
\label{eq:app-water-filling}
\end{equation}

These conditions characterize the global optimum because the allocation problem is
convex.

Features with larger \(\alpha_i\) generally have larger initial utility and therefore
receive more dimension. As their allocations increase, the compressed prediction
moves closer to the target and their marginal utilities fall. At the optimum, the
next infinitesimal unit of dimensional capacity has the same value \(\delta\)
wherever it is used.

Because all outputs share the same softmax denominator,
\[
u_i=u_i(\mathbf d)
\]
generally depends on the complete allocation, not only on \(d_i\). We therefore solve
Equation~\eqref{eq:app-allocation-program} directly rather than assuming separable
one-dimensional utility curves.

\subsection{Tight-frame realization and forced overlap}

Given \(\mathbf d^\star\), choose vectors \(E_i\in\mathbb R^d\) satisfying
\begin{equation}
\|E_i\|^2=d_i^\star,
\qquad
\sum_iE_iE_i^\top=I_d.
\label{eq:app-frame-realization}
\end{equation}
The Schur--Horn theorem guarantees that such a frame exists whenever
\[
0\le d_i^\star\le1,
\qquad
\sum_i d_i^\star=d.
\]

For any such realization,
\[
\sum_j
\langle E_i,E_j\rangle^2
=
E_i^\top
\left(
\sum_jE_jE_j^\top
\right)
E_i
=
\|E_i\|^2
=
d_i^\star.
\]
The term with \(j=i\) is
\[
\langle E_i,E_i\rangle^2
=
(d_i^\star)^2.
\]
Therefore
\begin{equation}
\boxed{
\sum_{j\neq i}
\langle E_i,E_j\rangle^2
=
d_i^\star(1-d_i^\star).
}
\label{eq:app-forced-overlap}
\end{equation}

Thus
\[
d_i^\star=0
\]
means feature \(i\) is absent,
\[
d_i^\star=1
\]
permits a dedicated direction, and
\[
0<d_i^\star<1
\]
requires nonzero overlap with other feature vectors.

The allocation fixes the total squared overlap associated with each feature, but
not its individual overlap partners or signs. Different tight-frame realizations
with the same row norms may therefore correspond to different orientations of the
same predicted allocation.

\subsection{The fixed-geometry decoder problem}

For a chosen tight-frame realization \(E^{(1)}\), define the predicted first-order
decoder and mean-embedding intercept by
\begin{equation}
\boxed{
\left(D^\star,e_0^\star\right)(E^{(1)})
=
\arg\min_{D,\,e_0}
\mathbb E_z
\operatorname{CE}
\left[
p_\star(\cdot\mid z),
\operatorname{softmax}
\left(
\left(z^\top E^{(1)}+e_0^\top\right)D
\right)
\right].
}
\label{eq:app-decoder-program}
\end{equation}
The intercept represents \(\bar E\), the mean of the state embedding: its logit
contribution \(e_0^\top D\) is a state-independent row through the decoder, so the
class stays bias-free. We fix the output gauge by requiring
\begin{equation}
D\mathbf1=0.
\label{eq:app-decoder-gauge}
\end{equation}
Subtracting the mean decoder vector changes every state's logits only by a constant
shared across outputs, and therefore does not change the predicted probabilities.

For fixed \(E^{(1)}\), the logits are linear in \(D\) at fixed \(e_0\) and linear in
\(e_0\) at fixed \(D\), so Equation~\eqref{eq:app-decoder-program} is convex in each
block; the pair is solved jointly.

The predicted first-order map is
\begin{equation}
\boxed{
L^{(1)}
=
E^{(1)}D^\star(E^{(1)}).
}
\label{eq:app-predicted-map}
\end{equation}
It automatically satisfies
\[
\operatorname{rank}L^{(1)}\le d.
\]

Let
\[
q(\cdot\mid z)
=
\operatorname{softmax}
\left(
\left(z^\top E^{(1)}+e_0^\top\right)D
\right),
\qquad
\bar r=\mathbb E_z\left[q-p_\star\right].
\]
The gradient with respect to \(D\) is
\[
\nabla_D\mathcal C
=
\mathbb E_z
\left[
\left(E^{(1)\top}z+e_0\right)
\left(q-p_\star\right)^\top
\right],
\]
and with respect to \(e_0\),
\[
\nabla_{e_0}\mathcal C
=
D\,\bar r.
\]
Therefore the joint optimum satisfies
\begin{equation}
\boxed{
E^{(1)\top}
\mathbb E_z
\left[
z\left(q-p_\star\right)^\top
\right]
=
-\,e_0\,\bar r^\top,
\qquad
D\,\bar r=0.
}
\label{eq:app-moment-matching}
\end{equation}

Thus the feature--output probability moments expressible through the chosen
\(d\)-dimensional geometry are matched up to a rank-one correction along the mean
residual \(\bar r\), and the mean residual lies in the null space of the decoder —
in particular the predicted map annihilates it, \(L^{(1)}\bar r=0\). These two
conditions are checkable observables of any claimed solution.

Because the allocation does not select a unique tight-frame orientation, the complete
prediction is the family
\[
\left\{
E^{(1)}D^\star(E^{(1)})
:
E^{(1)\top}E^{(1)}=I_d,\;
\|E_i\|^2=d_i^\star
\right\}.
\]

\subsection{Gauge and measurement}
\label{app:balanced}

The allocation \(d_i^\star\) and the first-order map \(E^{(1)}D\) are invariant to
hidden-coordinate reparameterization \(E\to EA\), \(D\to A^{-1}D\); raw feature and
decoder norms are not, and the loss cannot select them. Under gradient flow from small
initialization the two factors of a product \(ED^\top\) stay balanced,
\(E^\top E=D^\top D\), because the two gradient terms in \(\tfrac{d}{dt}(E^\top E-D^\top D)\)
cancel. Adam's per-coordinate rescaling breaks the cancellation: each row of the state
table receives the gradient of one state while each row of \(D\) sums over all states,
yet both move about one learning rate per step, so \(E\) grows roughly
\(\sqrt{n/m}\) times faster (measured 4.6--7.8 at \(n/m=32\); an SGD retrain is balanced to
\(0.4\%\)). The predicted model is stored in its own gauge too (\(E^\star\) Parseval, the
scale in \(D\)).

Every model's first-order vectors are the least-squares coefficients of its saved
representation on \([1,z]\) (the conditional-mean difference
\(\mathbb E[h\mid z_i=1]-\mathbb E[h\mid z_i=0]\) agrees in population but leaks the other
features' chance sample correlations, about \(1/\sqrt n\) each). Two readings follow.

\emph{Gauge-free.} The whitened frame \(\tilde E=E^{(1)}(E^{(1)\top}E^{(1)})^{-1/2}\) is
independent of \(A\) up to a rotation of \(\mathbb R^d\), so its Gram, the projector
\(E^{(1)}(E^{(1)\top}E^{(1)})^{+}E^{(1)\top}\), is invariant outright: its diagonal is
the leverage \(d_i\), its off-diagonals the overlaps of the frame. Allocation and every
Gram are read this way, on the feature side and on the readout side
(\(D^\top(DD^\top)^{+}D\)). To compare orientations across seeds, each seed's whitened
frame is rotated onto the first predicted frame by orthogonal Procrustes over all rows
in \(\mathbb R^d\) before averaging, and the Gram of the mean frame is shown: the mean of
per-seed Grams retains each seed's private wobble as haze, whereas the Gram of the mean
frame keeps the geometry the seeds share.

\emph{Balanced.} Sizes need a scale. Each model's row-centred logit table over states,
\(\Lambda=U\Sigma V^\top\) truncated to rank \(d\), is split evenly,
\begin{equation}
\boxed{
E_{\mathrm{bal}}
=
N^{1/4}\,U\Sigma^{1/2},
\qquad
D_{\mathrm{bal}}
=
N^{-1/4}\,\Sigma^{1/2}V^\top,
}
\label{eq:app-balanced-factorization}
\end{equation}
the same operation the other settings apply to their log-conditional tables, on the
per-state scale \(E_{\mathrm{bal}}^\top E_{\mathrm{bal}}/N=\operatorname{diag}(\sigma_k/\sqrt N)\).
Row-centring the table removes the softmax shift (a common vector added to every
readout row, \(0.8\) of a row's norm when trained) that would otherwise soak up leverage
on the readout side. The balancing acts on the table over states rather than on the
\(m\times m\) first-order map because the trained model's parameters are that table.
On this scale the Bayes table itself has sizes about \(\alpha_i/\sqrt{p(1-p)}\).

\subsection{Why strict additivity disfavors fractional allocation}

The allocation program predicts fractional values
\[
0<d_i^\star<1.
\]
Equation~\eqref{eq:app-forced-overlap} shows that any Euclidean realization of such
an allocation necessarily has overlapping first-order feature directions.

Suppose the complete representation were strictly additive:
\begin{equation}
E(z)
=
E_0+\sum_i z_iE_i.
\label{eq:app-strictly-additive}
\end{equation}
The complete logits would then be determined by the first-order map:
\[
\ell(z)
=
E_0^\top D
+
z^\top L^{(1)}.
\]
There would be no additional state dependence available to correct errors created
by overlap among the feature directions.

For small deviations from the ideal fractional prediction, the leading additive
realization cost is quadratic in the unwanted overlaps. If these overlaps have
approximately a common predictive price \(\beta>0\), then
\begin{equation}
\mathcal C_{\mathrm{add}}(\mathbf d)
\approx
\mathcal C(\mathbf d)
+
\beta
\sum_{i\neq j}
\langle E_i,E_j\rangle^2.
\label{eq:app-additive-overlap-cost}
\end{equation}
Using Equation~\eqref{eq:app-forced-overlap},
\begin{equation}
\boxed{
\mathcal C_{\mathrm{add}}(\mathbf d)
\approx
\mathcal C(\mathbf d)
+
\beta\sum_i d_i(1-d_i).
}
\label{eq:app-additive-allocation-cost}
\end{equation}

Since
\[
\sum_i d_i=d,
\]
we have
\[
\sum_i d_i(1-d_i)
=
d-\sum_i d_i^2.
\]
Thus the additive realization cost is minimized by making the allocation as
polarized as possible: some features approach \(d_i=1\), while the rest approach
\(d_i=0\).

The corresponding interior stationarity condition is
\begin{equation}
\boxed{
u_i(\mathbf d)
-
\beta(1-2d_i)
=
\delta.
}
\label{eq:app-additive-kkt}
\end{equation}
Relative to ordinary water filling,
\[
u_i(\mathbf d)=\delta,
\]
the extra term discourages small fractional allocations and favors completing
features that already receive substantial capacity.

Equation~\eqref{eq:app-additive-allocation-cost} is an approximation, not an exact
identity. The exact result is the forced-overlap law in
Equation~\eqref{eq:app-forced-overlap}. The approximation assumes that, in the iid
setting, the leading predictive cost of those overlaps is approximately proportional
to their total squared magnitude. The realization gap per unit of forced overlap measures \(\beta\)
(\(\approx0.02\)--\(0.04\) at the reference cell). Adding it to the allocation does
not improve the prediction there, because the completion refunds it; it is the
missing term where the pairs do not refund (steep ladders, correlated features), see
the refund assumption under Status.

\subsection{Why the full model can sustain fractional allocation}

The complete state embedding need not be strictly additive. In general,
\begin{equation}
E(z)
=
\bar E
+
\sum_i(z_i-p)E_i
+
E_{\ge2}(z).
\label{eq:app-full-representation}
\end{equation}
The residual \(E_{\ge2}(z)\) remains in the same space:
\[
E_{\ge2}(z)\in\mathbb R^d.
\]
It therefore introduces no additional representational dimensions.

Its role is instead to let different feature combinations occupy different points
inside the same \(d\)-dimensional space. For example, with two binary features,
strict additivity would require
\[
E(1,1)
=
E(1,0)+E(0,1)-E(0,0).
\]
A nonlinear state embedding may instead use
\[
E(1,1)
=
E(1,0)+E(0,1)-E(0,0)
+
\Delta_{12},
\]
where
\[
\Delta_{12}\in\mathbb R^d.
\]

Such context-dependent movements can correct predictive errors caused by overlapping
first-order directions. They therefore reduce the realization penalty associated
with fractional \(d_i\), allowing the first-order geometry of the full model to
remain close to the convex allocation prediction.

The main theory does not attempt to predict \(E_{\ge2}(z)\). Its role is to explain
why the convex first-order allocation can describe a full nonlinear model even when
a strictly additive rank-\(d\) representation prefers a much more polarized geometry.
The pair block of this correction is quantified in Section~\ref{sec:repair}.

\subsection{Summary}

The theory requires two convex optimizations.

First, the dimensional allocation is predicted by
\[
\boxed{
\mathbf d^\star
=
\arg\min_{\substack{0\le d_i\le1\\\sum_i d_i=d}}
\mathcal C(\mathbf d).
}
\]
Its KKT conditions give the resource-rational water-filling law
\[
u_i(\mathbf d^\star)=\delta
\]
for fractionally represented features.

Second, after choosing any tight-frame realization
\[
E^{(1)\top}E^{(1)}=I_d,
\qquad
\|E_i\|^2=d_i^\star,
\]
the first-order map is predicted by
\[
\boxed{
\left(D^\star,e_0^\star\right)
=
\arg\min_{D,\,e_0}
\mathbb E_z
\operatorname{CE}
\left[
p_\star(\cdot\mid z),
\operatorname{softmax}
\left(
\left(z^\top E^{(1)}+e_0^\top\right)D
\right)
\right].
}
\]
Then
\[
L^{(1)}=E^{(1)}D^\star,
\]
with the state-independent logit row \(e_0^{\star\top}D^\star\) carried through the
decoder.

The allocation predicts the whitened feature norms
\[
\|E_i\|^2=d_i^\star
\]
and the unavoidable squared overlap
\[
\sum_{j\neq i}\langle E_i,E_j\rangle^2
=
d_i^\star(1-d_i^\star).
\]

Finally, the whitened frame of the completed model gives the predicted allocation and
Grams, and the balanced rank-\(d\) factorization of its logit table over states gives
its sizes on the common scale (Section~\ref{app:balanced}).